% Lezione: 2026-10-01.
% Fonti: ../elaborati-assistente/trascrizioni-ricostruite/2026-10-01.txt;
% trascrizioni ricostruite per parte:
% ../elaborati-assistente/trascrizioni-ricostruite/2026-10-01_1.txt,
% ../elaborati-assistente/trascrizioni-ricostruite/2026-10-01_2.txt,
% ../elaborati-assistente/trascrizioni-ricostruite/2026-10-01_3.txt,
% ../elaborati-assistente/trascrizioni-ricostruite/2026-10-01_4.txt,
% ../elaborati-assistente/trascrizioni-ricostruite/2026-10-01_5.txt;
% trascrizioni grezze:
% ../trascarizioni/2026-10-01_1.txt, ../trascarizioni/2026-10-01_2.txt,
% ../trascarizioni/2026-10-01_3.txt, ../trascarizioni/2026-10-01_4.txt,
% ../trascarizioni/2026-10-01_5.txt;
% Le registrazioni corrispondenti sono risorse dell'utente; non sono fonti consultate.
% Le cinque parti sono incluse; non risultano parti mancanti. Alcuni passaggi testuali restano incerti.

\chapter{Introduzione alla fisiologia e al sistema nervoso}

\section{Fisiologia, ambiente e omeostasi}

Il corso si propone di studiare il sistema nervoso e la regolazione neuroendocrina. La fisiologia considera la struttura, la funzione e l'integrazione delle funzioni dell'organismo. Per comprendere il rapporto tra organismo e contesto è utile distinguere l'ambiente esterno dall'ambiente interno. La trascrizione descrive quest'ultimo come comprensivo degli spazi extracellulari esterni agli organi, ma la formulazione è incerta e non consente una precisazione ulteriore.

È importante distinguere la sensazione, relativa a ciò che viene rilevato dai sensi e dai recettori, dalla percezione, che comporta l'elaborazione centrale delle informazioni. I recettori possono far parte di organi di senso, come occhi, orecchie e lingua, essere proteine organizzate sulla membrana cellulare oppure corrispondere a singole cellule. Attraverso vie di comunicazione specifiche, le informazioni raggiungono centri di elaborazione. I centri di integrazione del sistema nervoso sono organizzati gerarchicamente e comunicano la risposta attraverso ulteriori vie, dirette agli effettori.

La risposta contribuisce al mantenimento dell'omeostasi e delle condizioni di salute. Il docente precisa che la fisiologia non va intesa in senso finalistico. Richiama inoltre la fitness collettiva e la riproduzione, ma il nesso tra i due concetti è poco chiaro nella trascrizione e non è possibile ricostruirlo con certezza.

\section{Sistema nervoso, percezione e comportamento}

Il sistema nervoso si trova all'interfaccia tra ambiente esterno e ambiente interno. L'interazione con il mondo permette apprendimento e memoria e comprende percezione sensoriale e motricità: la sensazione diventa percezione e l'organismo produce una risposta motoria. La trascrizione menziona, come esempio, la percezione delle proprie condizioni durante una lezione e la risposta dell'organismo, ma il passaggio specifico è poco intelligibile.

L'interazione avviene in stati di coscienza differenti. Il sonno non corrisponde allo spegnimento del sistema nervoso: durante il sonno il sistema continua a ricevere informazioni, comprese quelle provenienti dai propriocettori. La lezione richiama anche emozioni e comportamento. La rabbia è proposta come esempio di risposta che può concorrere alla difesa del territorio, della prole o del partner.

La trattazione del corso si concentra sulle basi fisiologiche del sistema nervoso. Le questioni filosofiche sul rapporto tra mente, cervello e anima sono esplicitamente al di fuori del suo ambito.

\section{Neuroni e cellule gliali}

Per affrontare un sistema complesso si adotta un approccio riduzionistico: si studiano prima componenti e meccanismi più semplici, per poi considerarne l'integrazione. I neuroni formano reti che operano in parallelo. Possono essere classificati in base alla morfologia, alla funzione e al neurotrasmettitore utilizzato. Tra gli esempi menzionati figurano neuroni bipolari e multipolari, neuroni sensoriali, motoneuroni e interneuroni, oltre a neuroni colinergici, GABAergici, serotoninergici, istaminergici e glutammatergici. Alcuni esempi e termini della classificazione sono corrotti nella trascrizione e non sono ricostruibili.

Accanto ai neuroni, il sistema nervoso comprende cellule gliali: sono citati astrociti, oligodendrociti, cellule di Schwann e microglia. La spiegazione delle funzioni dei singoli tipi cellulari è interrotta e in parte incomprensibile. La trascrizione riporta anche un confronto numerico tra glia e neuroni, ma il valore è ambiguo e non va corretto per supposizione.

\section{Potenziale di membrana e potenziale d'azione}

Il potenziale di membrana è una proprietà di tutte le cellule e riflette la distribuzione degli ioni e la permeabilità della membrana. Tra i fattori richiamati a lezione figurano le concentrazioni ioniche, la permeabilità relativa ai diversi ioni, la pompa sodio-potassio e la presenza intracellulare di proteine cariche non diffusibili. Le proprietà capacitive della membrana concorrono alle risposte passive; in alcuni tessuti si producono anche risposte attive.

Il potenziale d'azione è descritto come la registrazione della variazione del potenziale di membrana in funzione del tempo. La lezione richiama i potenziali di equilibrio degli ioni e i circuiti elettrici equivalenti della membrana; il nome dell'equazione menzionata non è comprensibile nella trascrizione. Il potenziale di membrana è una caratteristica di tutte le cellule. Le risposte passive dipendono dalle proprietà capacitive della membrana e dalle correnti che la attraversano; le risposte attive sono associate ai tessuti nervoso e muscolare, oltre a un terzo tessuto che non è distinguibile con certezza. Il potenziale dipende dalle concentrazioni ioniche, dalla permeabilità relativa della membrana ai diversi ioni, dalla pompa sodio-potassio e dalla presenza intracellulare di proteine cariche non diffusibili.

Il potenziale d'azione si genera nella regione iniziale dell'assone, dove convergono correnti ioniche associate alla sommazione spaziale e temporale delle risposte sinaptiche. In base alla legge di Ohm, le correnti producono variazioni del potenziale di membrana; queste possono modificare la conduttanza attraverso l'apertura dei canali ionici. Al raggiungimento della soglia, la corrente entrante di sodio contribuisce all'innesco del potenziale d'azione. La sommazione delle correnti provenienti dalle numerose sinapsi determina le variazioni di potenziale che possono portare all'innesco. Il valore numerico della soglia è ambiguo nella trascrizione e non viene riportato. Il docente invita a ripassare la fisiologia generale del potenziale di membrana e delle correnti.

\section{Canali ionici e proprietà del gating}

I canali ionici sono proteine integrali di membrana che consentono il passaggio selettivo di ioni, altrimenti ostacolato dal doppio strato lipidico. Il passaggio è rapido e genera una corrente che modifica il potenziale di membrana. Il repertorio dei canali varia tra tipi cellulari: anche cellule nervose diverse non esprimono necessariamente lo stesso insieme di canali. I canali presentano un meccanismo di apertura e chiusura, detto gating, e possono essere regolati dal voltaggio, da ligandi o da stimoli meccanici. La selettività fa sì che ciascun canale favorisca il passaggio di determinati ioni. La spiegazione del filtro di selettività si interrompe nella trascrizione e i dettagli della metafora enzimatica non sono ricostruibili.

Il gating descrive la transizione del canale tra stati funzionali, in particolare aperto e chiuso. Nelle registrazioni di singolo canale, la probabilità di apertura esprime la frazione di tempo in cui il canale è aperto. Per un canale regolato dal voltaggio, la relazione tra potenziale di membrana e probabilità di apertura permette di descriverne la sensibilità. Il potenziale di mezza attivazione corrisponde al valore di potenziale associato al 50\% di probabilità di apertura; la pendenza della curva descrive quanto rapidamente cambia la probabilità al variare del potenziale. La lezione collega queste proprietà alla diversa risposta delle popolazioni di canali e, di conseguenza, alle differenti modalità di scarica neuronale.

Nei canali voltaggio-dipendenti, il cambiamento conformazionale dei sensori precede l'apertura del poro. In un modello discusso a lezione, il canale si apre quando tutti e quattro i sensori si sono spostati; se se ne sposta solo una parte, il canale resta chiuso. Poiché questi eventi non avvengono necessariamente nello stesso momento nell'intera popolazione, l'apertura osservata è graduale e statistica. Il tempo necessario all'apertura e l'inattivazione sono richiamati in relazione ai canali del sodio.

Il periodo refrattario è il tempo durante il quale una cellula eccitabile non può rispondere a un nuovo stimolo oppure può farlo solo in risposta a uno stimolo più intenso. Nel periodo refrattario assoluto la cellula è insensibile; in quello relativo occorre uno stimolo maggiore. La refrattarietà è collegata alla disponibilità dei canali, che possono trovarsi nello stato chiuso, aperto o inattivato. Durante il recupero dall'inattivazione torna disponibile solo una parte della popolazione e la risposta a un secondo stimolo dipende dal tempo trascorso. I valori numerici esemplificativi sono confusi e non vengono specificati.

Le differenze tra canali nel potenziale di mezza attivazione e nella pendenza delle curve di attivazione contribuiscono a differenziare la risposta dei neuroni. Il docente paragona le diverse modalità di scarica a cadenze regionali differenti: il messaggio può essere riconoscibile pur variando la modalità con cui è espresso. Il dettaglio della codifica neuronale viene rimandato a lezioni successive.

Per i canali regolati da ligando, la probabilità di apertura dipende dalla concentrazione del ligando e dalla stechiometria dell'interazione; sono menzionati rapporti di legame diversi, senza che i relativi valori siano chiari. La corrente di un singolo canale dipende dalle proprietà del canale e dalle condizioni ioniche ai due lati della membrana. La lezione richiama in questo contesto la legge di Ohm e il ruolo della membrana come condensatore.

Il passaggio degli ioni attraverso il poro comporta la perdita del guscio di idratazione. La struttura interna del canale, comprese le catene laterali degli amminoacidi, contribuisce alla selettività e alla permeazione. Tossine e farmaci possono modificare le proprietà dei canali; il docente richiama anche l'importanza dell'azione transitoria di alcuni anestetici, ma la frase sui meccanismi e sui bersagli non è intelligibile. La trascrizione riporta circa 400 geni per canali ionici su un totale di circa 20.000 geni: questi numeri sono riportati come non verificati, non come dati certi.

I canali possono essere classificati in famiglie, tra cui quelle dei canali del sodio, del calcio e del potassio, ciascuna comprendente tipi con proprietà differenti. Il canale «funny» (If) è presentato come esempio presente nelle cellule pacemaker cardiache e coinvolto nella regolazione della frequenza cardiaca; il resto dell'esempio è interrotto da commenti in aula. La lezione collega la modulazione di If da parte delle catecolamine, attraverso recettori metabotropici, all'aumento della velocità di apertura del canale e della frequenza cardiaca.

\section{Neurotrasmettitori e recettori}

I recettori postsinaptici possono essere ionotropici, quando il recettore coincide con un canale, oppure metabotropici, quando il legame del trasmettitore avvia una cascata di segnalazione che può modulare altri canali. La lezione presenta i sistemi diffusi come circuiti di neuroni che utilizzano lo stesso trasmettitore e partecipano a funzioni specifiche. Per il sistema colinergico sono menzionati gruppi neuronali del prosencefalo basale e del ponte, con proiezioni verso la corteccia e altre regioni; alcune denominazioni anatomiche e destinazioni sono corrotte e non sono ricostruibili.

\subsection{Acetilcolina}

I recettori dell'acetilcolina comprendono recettori nicotinici ionotropici e recettori muscarinici metabotropici. Le denominazioni derivano rispettivamente da nicotina e muscarina. La muscarina viene descritta come componente tossica dell'Amanita muscaria, il fungo dal cappello rosso con macchie bianche; la farmacologia approfondirà questi recettori.

\subsection{Ammine biogene}

Tra le ammine biogene citate figurano dopamina, adrenalina, istamina, serotonina e noradrenalina. L'istamina viene menzionata in relazione ai mastociti e alle reazioni allergiche; la lezione ricorda anche neuroni istaminergici e cellule postsinaptiche con relativi recettori. Il rapporto tra sistema immunitario e sistema nervoso è riconosciuto, ma non approfondito.

La dopamina deriva dalla tirosina. Il sistema dopaminergico diffuso comprende la substantia nigra e l'area tegmentale ventrale; sono inoltre citati il sistema limbico e il nucleo accumbens. La degenerazione dei neuroni della substantia nigra è collegata al Parkinson e ai disturbi motori. La dopamina regola anche il rilascio di prolattina ed è coinvolta nei circuiti della ricompensa. La trascrizione menziona una terza regione anatomica, ma il nome è incerto.

Adrenalina e noradrenalina derivano dalla dopamina. Sono descritte come neurotrasmettitori del sistema nervoso periferico e possono essere rilasciate anche dalla midollare del surrene. Il sistema noradrenergico è associato al locus coeruleus, al talamo e al sistema limbico; non tutte le regioni elencate sono ricostruibili. Il docente collega questo sistema alla regolazione della soglia sonno-veglia. I recettori adrenergici alfa e beta sono metabotropici; i beta-bloccanti sono citati come esempio di antagonismo adrenergico e riduzione della frequenza cardiaca.

La serotonina deriva dal triptofano. La lezione la mette in relazione con la melatonina prodotta nell'epifisi e precisa che la serotonina cerebrale è sintetizzata localmente, poiché non attraversa la barriera ematoencefalica. Il sistema serotoninergico origina dai nuclei del rafe e proietta verso diverse regioni, incluso il midollo spinale. È associato alla regolazione del comportamento alimentare e sessuale e ai circuiti della ricompensa. L'esempio alimentare del docente non è sufficientemente chiaro per identificare l'alimento. I recettori serotoninergici possono essere ionotropici o metabotropici.

\subsection{Glutammato, GABA e altri mediatori}

Il glutammato è presentato come principale neurotrasmettitore eccitatorio. Tra i suoi recettori ionotropici sono citati AMPA, kainato e NMDA; sono menzionati anche recettori metabotropici, con effetti che possono essere eccitatori o inibitori. Il canale NMDA è descritto come bloccato dal magnesio a potenziali di riposo. Il glutammato attiva AMPA, la cui apertura depolarizza la membrana; la depolarizzazione rimuove il blocco del magnesio e permette l'attivazione del recettore NMDA. Quest'ultimo è descritto come permeabile soprattutto al calcio, oltre che ad altri ioni, e regolato sia dal ligando sia dal voltaggio.

La trascrizione attribuisce maggiore affinità a NMDA rispetto ad AMPA e kainato. Un'eccessiva attivazione glutammatergica e l'ingresso di calcio sono collegati a processi di eccitotossicità e danno cellulare, con attivazione di enzimi quali fosfolipasi, endonucleasi e proteasi. Gli astrociti partecipano alla rimozione del glutammato e alla protezione dall'eccitotossicità. Alcune spiegazioni molecolari sono disturbate dal rumore e non sono ricostruibili.

Il GABA è descritto come principale trasmettitore inibitorio del sistema nervoso centrale. I recettori GABA-A sono canali per il cloro; l'ingresso di cloro viene associato all'iperpolarizzazione. I recettori metabotropici GABA-B possono favorire l'apertura di canali del potassio; anche l'uscita di potassio è collegata all'iperpolarizzazione. La glicina è citata tra i neurotrasmettitori amminoacidici.

Il recettore metabotropico mGluR6 è menzionato in relazione ai circuiti retinici. L'ATP è collegato alla trasduzione del gusto: le cellule gustative lo rilasciano e attivano recettori ionotropici o metabotropici. Sono inoltre menzionati peptidi oppioidi endogeni, tra cui endorfine ed encefaline, altri piccoli peptidi e un possibile contributo all'asse intestino-cervello. Non tutti i peptidi elencati sono identificabili con sicurezza. Infine, sono citati gas come possibili neurotrasmettitori, tra cui ossido nitrico e acido solfidrico, ed endocannabinoidi.

In chiusura, il docente ricorda che la lezione contiene domande d'esame e invita a ripassare potenziale di membrana, sinapsi, potenziamenti e correnti. L'ultima frase è in parte incomprensibile.

\section{Sintesi}

La lezione introduce il sistema nervoso come rete integrativa che mette in relazione ambiente, percezione, risposta e omeostasi. A livello cellulare, le proprietà della membrana e dei canali ionici permettono la generazione e la modulazione dei segnali elettrici. La comunicazione sinaptica dipende da neurotrasmettitori e recettori con proprietà differenti. Gli argomenti e i dettagli sostanziali ricostruibili dalle trascrizioni delle cinque parti sono riportati nel capitolo; i passaggi illeggibili o incerti restano qualificati come tali e non sono completati con informazioni esterne.
